%% Annexe E : Bibliographie commentée
\chapter{Bibliographie Commentée}
\label{ann:bibliographie}

Pour chaque partie, quelques lectures. Les références complètes figurent dans
la bibliographie du livre.

\section*{Parties I et II · Superviser, collecter, détecter}
\begin{itemize}
    \item NIST SP 800-61r2~\cite{nist80061r2}~: le traitement des incidents.
    \item CIS Controls v8~\cite{cis2021controls}~: la place de la supervision.
    \item MITRE ATT\&CK~\cite{mitreattack}~: le vocabulaire des techniques.
\end{itemize}

\section*{Partie III · Cadre normatif et juridique}
\begin{itemize}
    \item ISO/IEC 27035~\cite{iso27035} et 27037~\cite{iso27037}~: incident et preuve.
    \item NIST SP 800-86~\cite{nist80086}, RFC 3227~\cite{rfc3227}~: forensique et volatilité.
    \item Lois 2010/012~\cite{loi2010012}, 2024/017~\cite{loi2024017}~; Malabo~\cite{malabo2014}~; Budapest~\cite{budapest2001}.
\end{itemize}

\section*{Partie IV · L'analyse forensique}
\begin{itemize}
    \item NIST SP 800-86~\cite{nist80086}~: la démarche forensique.
    \item ISO/IEC 27037~\cite{iso27037}~: la préservation.
\end{itemize}

\section*{Partie V · La réponse à incident}
\begin{itemize}
    \item ISO/IEC 27035~\cite{iso27035}~: le cycle de réponse.
    \item Chaîne d'attaque~\cite{killchain2011}, Diamant~\cite{diamond2013}, STIX/TAXII~\cite{stixtaxii}.
\end{itemize}

\section*{Partie VI · Apprendre et mesurer}
\begin{itemize}
    \item ISO/IEC 27004~\cite{iso27004}~: la mesure.
    \item T.~E.~Minka Mi Nguidjoï~\cite{minka2026auditologie}~: le verdict à trois valeurs.
\end{itemize}
